\documentclass[
  12pt,
  oneside,
  a4paper,
  brazil
]{abntex2}

\usepackage[T1]{fontenc}
\usepackage[utf8]{inputenc}
\usepackage{lmodern}
\usepackage{indentfirst}
\usepackage{microtype}
\usepackage{amsmath}
\usepackage{array}

\titulo{Técnicas de Otimização Aplicadas à Alocação de Recursos Financeiros em Ativos da B3}
\autor{André Sapucaia De Araujo}
\orientador{Rogério de Oliveira}
\local{São Paulo}
\data{2026}
\instituicao{%
  Universidade Presbiteriana Mackenzie
  \par
  Curso de Ciência da Computação}
\tipotrabalho{Trabalho de Conclusão de Curso}
\preambulo{Trabalho de Conclusão de Curso apresentado ao Curso de Ciência da Computação da Universidade Presbiteriana Mackenzie, como requisito parcial para obtenção do título de Bacharel em Ciência da Computação.}

\begin{document}

\imprimircapa

\begin{resumo}
A alocação eficiente de recursos financeiros constitui um dos principais desafios das finanças quantitativas modernas, especialmente diante da crescente complexidade dos mercados e da elevada disponibilidade de ativos negociados em bolsas de valores. Nesse contexto, técnicas de otimização vêm sendo amplamente utilizadas para auxiliar investidores e instituições financeiras na construção de carteiras de investimento capazes de maximizar retornos e minimizar riscos.

O presente trabalho tem como objetivo analisar diferentes técnicas de otimização aplicadas à alocação de recursos financeiros em ativos negociados na B3, considerando simultaneamente critérios de retorno esperado, risco, volatilidade e diversificação. A pesquisa fundamenta-se em abordagens clássicas e heurísticas de otimização, incluindo Programação Quadrática, Algoritmos Genéticos e Simulated Annealing.

A metodologia proposta baseia-se na coleta de dados históricos de ações brasileiras por meio da plataforma Yahoo Finance, contemplando ativos representativos do mercado financeiro nacional no período entre 2015 e 2025. Os dados serão tratados por meio do cálculo de retornos logarítmicos e matriz de covariância, permitindo a modelagem matemática do problema de otimização de portfólio.

Como resultados esperados, busca-se identificar vantagens e limitações das técnicas analisadas, bem como avaliar sua aplicabilidade no contexto do mercado financeiro brasileiro. Espera-se ainda demonstrar que métodos heurísticos podem oferecer maior flexibilidade em cenários complexos, enquanto abordagens clássicas apresentam maior estabilidade computacional.

\textbf{Palavras-chave}: otimização de portfólio; algoritmos genéticos; simulated annealing; mercado financeiro; finanças quantitativas.
\end{resumo}

\begin{resumo}[Abstract]
\begin{otherlanguage*}{english}
Efficient financial resource allocation represents one of the main challenges in modern quantitative finance, especially given the increasing complexity of financial markets and the large availability of assets traded on stock exchanges. In this context, optimization techniques have been widely applied to support investors and financial institutions in the construction of investment portfolios capable of maximizing returns while minimizing risks.

This study aims to analyze different optimization techniques applied to the allocation of financial resources in assets traded on B3, simultaneously considering expected return, risk, volatility, and diversification criteria. The research is based on classical and heuristic optimization approaches, including Quadratic Programming, Genetic Algorithms, and Simulated Annealing.

The proposed methodology is based on the collection of historical data from Brazilian stocks through the Yahoo Finance platform, considering representative assets from the Brazilian financial market between 2015 and 2025. The data will be processed through logarithmic return calculations and covariance matrix estimation, enabling the mathematical modeling of the portfolio optimization problem.

As expected results, the study seeks to identify advantages and limitations of the analyzed techniques, as well as evaluate their applicability within the Brazilian financial market context. It is also expected to demonstrate that heuristic methods may provide greater flexibility in complex scenarios, while classical approaches offer greater computational stability.

\textbf{Keywords}: portfolio optimization; genetic algorithms; simulated annealing; financial market; quantitative finance.
\end{otherlanguage*}
\end{resumo}

\textual

\chapter{Introdução}
O mercado financeiro moderno caracteriza-se por um ambiente altamente dinâmico, competitivo e sujeito a constantes oscilações econômicas. Nesse cenário, investidores e instituições financeiras buscam continuamente estratégias capazes de maximizar o retorno de seus investimentos ao mesmo tempo em que minimizam os riscos associados às incertezas do mercado. A crescente complexidade das operações financeiras e o aumento da disponibilidade de dados impulsionaram o desenvolvimento de técnicas quantitativas aplicadas à análise e otimização de carteiras de investimento.

A alocação eficiente de recursos financeiros constitui um dos principais problemas estudados no campo das finanças quantitativas. A construção de um portfólio envolve a distribuição de capital entre diferentes ativos financeiros, considerando fatores como retorno esperado, volatilidade, correlação entre ativos e tolerância ao risco do investidor. Dessa forma, a tomada de decisão relacionada à composição de carteiras torna-se um problema matemático complexo, especialmente em mercados com grande quantidade de ativos e variáveis envolvidas.

Segundo Zhang e Li (2020), a utilização de técnicas modernas de otimização permite explorar padrões complexos nos mercados financeiros, contribuindo para a obtenção de soluções mais eficientes na construção de portfólios. Além disso, Zhou, Li e Zhang (2022) destacam que algoritmos heurísticos apresentam elevado potencial em problemas financeiros caracterizados por múltiplas restrições e grande dimensionalidade.

Tradicionalmente, modelos de otimização financeira utilizam abordagens clássicas baseadas em programação matemática, como a Programação Quadrática. Entretanto, tais métodos podem apresentar limitações em cenários mais complexos, especialmente quando o problema envolve funções não lineares ou grande quantidade de restrições operacionais. Em razão disso, técnicas heurísticas e metaheurísticas passaram a ser amplamente exploradas no contexto financeiro.

Entre as principais abordagens heurísticas aplicadas à otimização de portfólios destacam-se os Algoritmos Genéticos e o Simulated Annealing. Essas técnicas possuem capacidade de explorar espaços de busca complexos e escapar de mínimos locais, tornando-se alternativas relevantes em problemas de otimização multiobjetivo.

Diante desse contexto, o presente trabalho propõe a análise de diferentes técnicas de otimização aplicadas à alocação de recursos financeiros em ativos negociados na B3, buscando comparar suas características, vantagens, limitações e aplicabilidade prática no mercado financeiro brasileiro.

O objetivo geral deste estudo consiste em analisar e comparar técnicas de otimização aplicadas à construção de carteiras de investimento utilizando ativos negociados na B3.

Como objetivos específicos, busca-se:
\begin{itemize}
  \item coletar e tratar dados históricos de ações brasileiras;
  \item modelar matematicamente o problema de otimização de portfólio;
  \item implementar algoritmos clássicos e heurísticos;
  \item comparar os métodos utilizando métricas financeiras;
  \item analisar o desempenho das carteiras geradas.
\end{itemize}

A relevância deste trabalho justifica-se pela crescente importância da análise quantitativa no mercado financeiro e pela necessidade de desenvolvimento de estratégias mais eficientes de alocação de recursos. Além disso, a pesquisa contribui para o entendimento da aplicação de técnicas computacionais modernas em problemas financeiros reais.

Este trabalho está estruturado em quatro capítulos principais. Inicialmente, apresenta-se o referencial teórico relacionado à teoria moderna do portfólio e às técnicas de otimização estudadas. Em seguida, descreve-se a metodologia adotada, incluindo coleta de dados, modelagem matemática e métricas de avaliação.

\chapter{Referencial Teórico}
\section{Mercado Financeiro e Carteiras de Investimento}
O mercado financeiro é responsável pela intermediação de recursos entre agentes superavitários e deficitários, viabilizando a formação de poupança, financiamento de empresas e alocação de capital em diferentes classes de ativos. No contexto acionário, a decisão de investimento envolve a escolha de ativos com perfis distintos de retorno, risco, liquidez e sensibilidade a fatores macroeconômicos.

No caso brasileiro, a B3 reúne empresas de diferentes setores e níveis de maturidade, criando um ambiente heterogêneo para o investidor. Essa heterogeneidade, embora amplie as oportunidades de diversificação, também aumenta a complexidade analítica do processo de seleção de ativos. Dessa forma, a simples escolha baseada em desempenho passado não é suficiente para compor carteiras robustas no longo prazo.

Uma carteira de investimentos pode ser definida como um vetor de pesos aplicado sobre um conjunto de ativos, no qual cada peso representa a proporção do capital alocada em um ativo específico. A qualidade dessa carteira depende da relação conjunta entre retorno esperado e risco assumido. Conforme discutem Gu, Kelly e Xiu (2020), decisões de alocação mais eficientes exigem tratamento estatístico adequado dos dados históricos e modelagem quantitativa consistente.

\subsection{Relação Risco-Retorno}
A literatura de finanças estabelece que retornos mais elevados geralmente estão associados a maior incerteza. Em termos operacionais, essa incerteza é frequentemente mensurada pela volatilidade dos retornos e pela covariância entre ativos. Assim, duas carteiras com o mesmo retorno esperado podem apresentar riscos diferentes, dependendo de como os ativos se relacionam entre si.

No problema de alocação, não basta observar o risco individual de cada ativo; é necessário considerar também o risco conjunto do portfólio. Esse ponto justifica o uso de modelos baseados em matriz de covariância, pois a contribuição de risco de cada ativo depende simultaneamente de seu peso e de sua correlação com os demais componentes da carteira.

\subsection{Diversificação e Eficiência}
A diversificação é um princípio central da gestão de carteiras. Ao combinar ativos com comportamentos parcialmente descorrelacionados, é possível reduzir o risco total sem, necessariamente, reduzir o retorno esperado na mesma proporção. Em termos matemáticos, essa propriedade aparece na estrutura quadrática do risco da carteira.

Contudo, diversificar não significa distribuir capital de forma uniforme em todos os ativos. A diversificação eficiente depende de otimização sob restrições, considerando limites de alocação, aversão ao risco e critérios de viabilidade prática. Por isso, técnicas de otimização são fundamentais para transformar o princípio de diversificação em estratégia aplicável.

\section{Teoria Moderna do Portfólio}
A Teoria Moderna do Portfólio (TMP), consolidada por Markowitz, formaliza a construção de carteiras eficientes por meio da otimização conjunta de retorno e risco. Embora o modelo clássico seja amplamente conhecido, sua lógica permanece atual em aplicações contemporâneas, inclusive quando combinada com métodos heurísticos.

Seja $w$ o vetor de pesos dos ativos, $\mu$ o vetor de retornos esperados e $\Sigma$ a matriz de covariância. O retorno esperado da carteira é dado por:
\[
E(R_p)=w^T\mu
\]

e a variância da carteira por:
\[
\sigma_p^2=w^T\Sigma w
\]

A partir dessas expressões, pode-se formular problemas de minimização de risco para um retorno alvo, ou de maximização de retorno para um nível de risco aceitável.

\subsection{Fronteira Eficiente}
A fronteira eficiente representa o conjunto de carteiras não dominadas no plano risco-retorno. Uma carteira é eficiente quando não existe outra carteira que ofereça maior retorno para o mesmo risco, ou menor risco para o mesmo retorno. Esse conceito é central para comparação entre estratégias de otimização.

Na prática, a fronteira eficiente serve como referência para avaliar se um algoritmo está gerando soluções competitivas. Métodos diferentes podem atingir regiões distintas da fronteira, com variações em estabilidade numérica, tempo de execução e aderência a restrições operacionais.

\subsection{Limitações do Modelo Clássico}
Apesar de sua relevância, o modelo clássico apresenta limitações importantes: sensibilidade a estimativas de retorno, dependência de hipóteses estatísticas fortes e dificuldade para lidar com restrições não lineares complexas. Essas limitações estimulam o uso de abordagens heurísticas e metaheurísticas, especialmente em cenários de alta dimensionalidade e incerteza elevada.

Kumar e Singh (2022) destacam que modelos multiobjetivo e técnicas evolutivas ampliam a capacidade de busca por soluções eficientes quando o espaço de decisão se torna mais complexo. Assim, a TMP permanece como base conceitual, mas frequentemente é complementada por algoritmos mais flexíveis.

\section{Técnicas de Otimização Aplicadas a Portfólios}
As técnicas de otimização em finanças podem ser divididas, de forma geral, em métodos determinísticos e métodos heurísticos. Métodos determinísticos tendem a ter maior previsibilidade e convergência controlada sob hipóteses bem definidas. Métodos heurísticos, por sua vez, priorizam flexibilidade de busca e adaptação a superfícies de solução irregulares.

\subsection{Programação Quadrática}
A Programação Quadrática (PQ) modela o problema de alocação considerando uma função objetivo quadrática e restrições lineares. A formulação típica de utilidade risco-retorno é:
\[
\max \left(\mu^T w - \lambda w^T\Sigma w\right)
\]

sujeita a:
\[
\sum_{i=1}^{n}w_i=1, \quad w_i\geq 0
\]

em que $\lambda$ representa o parâmetro de aversão ao risco. Quanto maior o valor de $\lambda$, maior a penalização da volatilidade da carteira.

Entre as vantagens da PQ destacam-se: eficiência computacional, interpretação econômica clara e boa estabilidade em problemas convexos. Entre as limitações, destacam-se menor flexibilidade para restrições não convencionais e sensibilidade a parâmetros estimados com ruído.

\subsection{Algoritmos Genéticos}
Os Algoritmos Genéticos (AG) são inspirados em processos evolutivos e operam sobre uma população de soluções candidatas. Em cada geração, indivíduos com melhor aptidão têm maior probabilidade de reprodução, enquanto operadores de cruzamento e mutação promovem exploração do espaço de busca.

No contexto de portfólios, cada indivíduo pode representar um vetor de pesos, normalizado para respeitar restrições de orçamento e não negatividade. A função de aptidão geralmente combina retorno esperado e penalização de risco, em linha com a função objetivo do problema.

De acordo com Bai, Liu e Wang (2020), AG apresenta bom desempenho em problemas com múltiplas restrições e superfícies de busca não lineares. Contudo, seu desempenho depende da calibração de hiperparâmetros, como tamanho da população, taxa de mutação e número de gerações.

\subsection{Simulated Annealing}
O Simulated Annealing (SA) é um método probabilístico baseado no processo de recozimento térmico. A estratégia consiste em aceitar, com certa probabilidade, soluções temporariamente piores para evitar aprisionamento em ótimos locais. Essa probabilidade decresce ao longo das iterações segundo um cronograma de resfriamento.

No problema de alocação de ativos, o SA parte de um vetor de pesos inicial e executa perturbações controladas, avaliando a qualidade de cada nova solução por meio da função objetivo risco-retorno. A aceitação de movimentos de piora no início da busca amplia exploração global; no fim da busca, a redução da temperatura favorece refinamento local.

Conforme Li e Wang (2020), o SA é especialmente útil quando o espaço de decisão contém múltiplas regiões de atratividade e descontinuidades práticas de modelagem. Sua principal desvantagem é o possível aumento de custo computacional, dependendo do número de iterações e da estratégia de resfriamento.

\section{Métricas Financeiras para Avaliação de Portfólios}
A comparação entre algoritmos exige um conjunto de métricas capazes de capturar retorno, risco e eficiência operacional. Neste trabalho, são consideradas as seguintes métricas principais:
\begin{itemize}
  \item \textbf{Retorno esperado}: média dos retornos estimados da carteira.
  \item \textbf{Volatilidade}: desvio padrão dos retornos da carteira, utilizado como proxy de risco.
  \item \textbf{índice de Sharpe}: razão entre o retorno esperado e a volatilidade (sem taxa livre de risco, calculada sobre retornos diários), indicando retorno ajustado ao risco.
  \item \textbf{Drawdown}: perda acumulada máxima a partir de um pico, útil para avaliação de risco de queda.
  \item \textbf{Tempo computacional}: custo de processamento para obtenção da solução.
\end{itemize}

Além dessas métricas, a análise de contribuição de risco por ativo permite interpretar como cada posição impacta o risco agregado da carteira. Essa interpretação é relevante para decisões de rebalanceamento e controle de concentração.

\section{Trabalhos Relacionados}
A literatura recente aponta crescimento consistente do uso de inteligência computacional em finanças quantitativas. Zhang e Li (2020) apresentam uma revisão sobre integração entre aprendizado de máquina e otimização de portfólios, destacando ganhos de desempenho em cenários complexos.

Heaton, Polson e Witte (2020) discutem aplicações de deep learning em construção de portfólios dinâmicos, evidenciando que modelos avançados podem capturar não linearidades relevantes da dinâmica de mercado. Em paralelo, Zhou, Li e Zhang (2022) mostram que metaheurísticas mantêm competitividade quando comparadas a métodos tradicionais sob múltiplas restrições.

Apesar dos avanços, ainda há lacunas para mercados emergentes, sobretudo quanto à validação empírica com ativos locais e restrições realistas de alocação. Nesse sentido, o mercado brasileiro oferece contexto apropriado para avaliação comparativa de técnicas clássicas e heurísticas.

\chapter{Metodologia}
\section{Caracterização da Pesquisa}
Esta pesquisa possui natureza aplicada, com abordagem quantitativa e objetivo comparativo-experimental. O foco está em avaliar o desempenho de três técnicas de otimização na construção de carteiras com ativos negociados na B3, utilizando critérios padronizados de entrada, função objetivo e métricas de saída.

A estratégia metodológica combina etapas de engenharia de dados, modelagem matemática, implementação computacional e análise de resultados. O desenho experimental foi estruturado para garantir reprodutibilidade, modularidade do código e comparabilidade entre algoritmos.

\section{Base de Dados e Critérios de Seleção}
Os dados históricos de preços são obtidos por meio da plataforma Yahoo Finance, com horizonte temporal de janeiro de 2015 a dezembro de 2025, o que resulta em 2.736 observações diárias de retorno. A seleção de ativos considera critérios de liquidez, representatividade setorial e presença recorrente entre papéis de maior negociação no mercado brasileiro. Nos experimentos apresentados, foram utilizados três ativos de grande liquidez: PETR4, VALE3 e ITUB4.

Os dados brutos são coletados em frequência diária e armazenados com mecanismo de cache local em formato CSV, permitindo fallback em situações de indisponibilidade temporária da API. Esse procedimento aumenta estabilidade operacional do experimento e reduz impacto de limitações de taxa de consulta.

\subsection{Ativos e Janela Temporal}
A janela temporal é definida para capturar diferentes regimes de mercado, incluindo períodos de maior e menor volatilidade. Isso permite analisar a robustez dos algoritmos em cenários variados, evitando conclusões baseadas apenas em intervalos de baixa variabilidade.

A composição final de ativos pode ser ajustada conforme objetivo da análise, preservando-se o mesmo protocolo de pré-processamento e comparação para todos os métodos.

\section{Pré-processamento dos Dados}
Após coleta, os dados passam por verificação de consistência, remoção de registros inválidos e alinhamento temporal entre ativos. Em seguida, são calculados retornos logarítmicos diários:
\[
r_t = \ln\left(\frac{P_t}{P_{t-1}}\right)
\]

A escolha por retornos logarítmicos se justifica por propriedades estatísticas úteis em modelagem financeira, como aditividade temporal aproximada e melhor comportamento em séries de preços multiplicativas.

\subsection{Estimativas Estatísticas}
Com base nos retornos, são estimados:
\begin{itemize}
  \item vetor de retornos esperados ($\mu$);
  \item vetor de volatilidades individuais;
  \item matriz de covariância ($\Sigma$).
\end{itemize}

Essas estimativas alimentam a função objetivo comum dos três algoritmos, garantindo comparabilidade metodológica. Sempre que necessário, são aplicados tratamentos numéricos para evitar instabilidades, como regularização de valores muito próximos de zero.

\section{Modelagem Matemática do Problema}
O problema de alocação é formulado como a maximização de uma utilidade média-variância, isto é, o retorno esperado penalizado pela variância da carteira:
\[
\max f(w)=\mu^T w-\lambda\, w^T\Sigma w
\]

sujeito a:
\[
\sum_{i=1}^{n}w_i=1, \quad w_i\geq 0, \quad w_i\leq w_{max}
\]

em que $w_{max}$ limita concentração excessiva em um único ativo e $\lambda$ controla aversão ao risco. Nos experimentos, adotaram-se $\lambda = 3$ e $w_{max} = 0{,}6$. Essa formulação é aplicada de forma idêntica para Programação Quadrática, Algoritmo Genético e Simulated Annealing.

\subsection{Padronização Experimental}
Para evitar viés de comparação, os algoritmos compartilham:
\begin{itemize}
  \item o mesmo conjunto de ativos e janela temporal;
  \item a mesma função objetivo;
  \item as mesmas restrições de alocação;
  \item o mesmo formato de saída com métricas padronizadas.
\end{itemize}

Essa padronização permite atribuir diferenças de desempenho ao comportamento do algoritmo, e não a diferenças de entrada ou de definição do problema.

\section{Implementação dos Algoritmos}
Foram implementados três métodos: Programação Quadrática, Algoritmo Genético e Simulated Annealing. A implementação foi realizada em Python, com foco em simplicidade, reprodutibilidade e execução totalmente local.

\subsection{Programação Quadrática}
A abordagem determinística busca soluções de alta estabilidade, respeitando restrições do problema e fornecendo baseline analítico para comparação com métodos heurísticos. O problema é resolvido pelo método SLSQP da biblioteca SciPy; como a função objetivo é côncava e as restrições são lineares, a solução encontrada é o ótimo global.

\subsection{Algoritmo Genético}
A abordagem evolutiva utiliza população inicial, seleção por aptidão, cruzamento e mutação, com normalização dos pesos a cada iteração para preservação das restrições de viabilidade. Foram utilizados população de 32 indivíduos, 50 gerações e semente aleatória 42.

\subsection{Simulated Annealing}
A abordagem probabilística executa perturbações graduais nos pesos e aceita soluções de piora de forma controlada por temperatura, favorecendo escape de ótimos locais e refinamento progressivo. Foram utilizadas 250 iterações, temperatura inicial 1,0, taxa de resfriamento 0,98 e semente aleatória 42.

\section{Métricas de Avaliação e Comparação}
A avaliação dos algoritmos utiliza métricas financeiras e computacionais:
\begin{itemize}
  \item valor da função objetivo;
  \item retorno esperado da carteira;
  \item volatilidade do portfólio;
  \item índice de Sharpe;
  \item drawdown máximo;
  \item tempo de execução em milissegundos.
\end{itemize}

Também é analisada a distribuição de pesos e a contribuição percentual de risco por ativo, para verificar concentração e coerência econômica das soluções.

\subsection{Ranking de Desempenho}
Com base nas métricas, é gerado um ranking comparativo entre os três métodos. O critério principal prioriza função objetivo e desempenho ajustado ao risco, sem desconsiderar custo computacional. Essa visão multicritério evita conclusões simplificadas baseadas em apenas uma métrica.

\section{Backtesting e Validação}
Além das métricas pontuais, realiza-se backtesting das carteiras para observar a evolução acumulada de desempenho e o comportamento em períodos de estresse. Nesta etapa, o backtest utiliza todo o período amostral (in-sample), isto é, os mesmos dados empregados na otimização; a validação fora da amostra, com divisão entre períodos de treino e de teste, está prevista como etapa seguinte do trabalho.

Para reforçar a validade dos resultados, recomenda-se execução com múltiplos conjuntos de ativos e análise de sensibilidade dos hiperparâmetros dos métodos heurísticos.

\section{Procedimentos de Reprodutibilidade}
A arquitetura do projeto foi dividida em módulos de ingestão, engenharia de atributos, algoritmos, avaliação e infraestrutura. Essa separação facilita auditoria, testes e evolução incremental do trabalho.

A execução local pode ser reproduzida por scripts padronizados e variáveis de ambiente. O código-fonte é versionado em repositório Git e todo o experimento é executado localmente, por meio de API e interface web próprias, sem dependência de serviços em nuvem.

Essa estratégia metodológica garante rastreabilidade entre dados de entrada, parâmetros de execução e resultados finais, aspecto essencial para robustez acadêmica do estudo.


\postextual
\chapter*{Referências}
\addcontentsline{toc}{chapter}{Referências}

BAI, X.; LIU, J.; WANG, Y. A hybrid genetic algorithm for portfolio optimization. \textit{Expert Systems with Applications}, v. 161, 2020.

GU, S.; KELLY, B.; XIU, D. Empirical asset pricing via machine learning. \textit{The Review of Financial Studies}, v. 33, n. 5, p. 2223--2273, 2020.

HEATON, J.; POLSON, N.; WITTE, J. Deep learning for finance: deep portfolios. \textit{Applied Stochastic Models in Business and Industry}, v. 36, n. 1, p. 3--12, 2020.

KUMAR, R.; SINGH, S. Multi-objective portfolio optimization using evolutionary algorithms. \textit{Computers \& Operations Research}, v. 140, 2022.

LI, Y.; WANG, S. Simulated annealing-based portfolio optimization under uncertainty. \textit{Applied Soft Computing}, v. 95, 2020.

ZHANG, Y.; LI, X. Portfolio optimization using machine learning: A review. \textit{IEEE Access}, v. 8, p. 203960--203976, 2020.

ZHOU, Y.; LI, H.; ZHANG, J. Metaheuristic approaches for portfolio optimization: A systematic review. \textit{Swarm and Evolutionary Computation}, v. 68, 2022.

\end{document}
